% programme_summary_sections5_8.tex
% Sections 5--8 of the programme summary paper
% To be merged with programme_summary_section1.tex and
% programme_summary_sections2_4.tex before \end{document}.

% ====================================================================
\section{The collision residue and the seven faces}
\label{sec:seven-faces}
% ====================================================================

The bar differential extracts OPE data at collision divisors.
When two points collide, the residue of the bar-intrinsic MC
element $\Theta_\cA$ at the codimension-one stratum
$D_{\{1,2\}} \subset \FM_2(X)$ defines the \emph{collision
residue}:
\begin{equation}\label{eq:collision-residue}
r_\cA(z)
\;:=\;
\Res^{\mathrm{coll}}_{0,2}(\Theta_\cA)
\;\in\;
\End_\cA(2)[\![z^{-1}]\!]\,.
\end{equation}
The collision residue is the arity-$2$, genus-$0$ projection
of $\Theta_\cA$. Its pole orders are one less than those of
the OPE: the logarithmic kernel $d\log(z-w)$ absorbs one
power of $(z - w)$. For a chiral algebra with OPE poles of
maximal order~$p$, the collision residue has poles of
order~$p - 1$.

% ====================================================================
\subsection{Pole-order dichotomy}
\label{ssec:pole-dichotomy}
% ====================================================================

The maximal pole order of the collision residue separates the
standard landscape into two regimes:

\begin{center}
\renewcommand{\arraystretch}{1.3}
\begin{tabular}{clcl}
\toprule
\textbf{Pole order} & \textbf{Algebra}
 & \textbf{Collision} & \textbf{Class} \\
\midrule
$1$ & Heisenberg $\cH_k$ & $k/z$ & $\mathbf{G}$ \\
$1$ & Kac--Moody $\widehat{\fg}_k$ & $\Omega/((k{+}h^\vee)z)$ & $\mathbf{L}$ \\
$3$ & Virasoro $\Vir_c$ & $(c/2)/z^3 + 2T/z$ & $\mathbf{M}$ \\
$2N{-}1$ & $\mathcal{W}_N$ & poles through $z^{-(2N-1)}$
 & $\mathbf{M}$ \\
\bottomrule
\end{tabular}
\end{center}

\noindent
This table encodes a structural dichotomy: algebras whose
collision residue has at most a simple pole (classes~$\mathbf{G}$
and~$\mathbf{L}$) have finite shadow depth and
Swiss-cheese-formal bar complexes; algebras with higher-order
poles (class~$\mathbf{M}$) have infinite shadow towers and
genuinely non-formal $A_\infty$-structure.

The dichotomy is not a property of individual OPE coefficients
but of the collision residue as a whole: the $d\log$ absorption
converts the OPE's double pole (carrying the metric) and simple
pole (carrying the Lie bracket) into a single simple pole
(carrying the Casimir tensor) for Kac--Moody, while the
Virasoro OPE's quartic pole becomes a cubic collision pole,
producing an $r$-matrix with genuine spectral dependence.

% ====================================================================
\subsection{Seven faces of $r(z)$}
\label{ssec:seven-faces}
% ====================================================================

The collision residue is a single object viewed through seven
algebraic lenses.

\begin{theorem}[Seven-face identification]
\label{thm:seven-faces}
For a modular Koszul chiral algebra~$\cA$, the following seven
expressions are equal as elements of\/
$\End_\cA(2)[\![z^{-1}]\!]$:
\begin{enumerate}[label=\textup{(F\arabic*)}]
\item The bar collision residue:
 $r_\cA(z)
 = \Res^{\mathrm{coll}}_{0,2}(\Theta_\cA)$.
\item The spectral $R$-matrix of the ordered bar:
 $R(z) = \mathrm{id} + r(z)/z + \cdots$ solves the
 classical Yang--Baxter equation.
\item The $\lambda$-bracket of the Poisson vertex algebra:
 $r(z) = \int_0^\infty e^{-\lambda z}
 \{a_\lambda b\}\,d\lambda$
 \textup{(}Laplace transform\textup{)}.
\item The Knizhnik--Zamolodchikov connection:
 $\nabla_{\mathrm{KZ}}
 = d - r(z)\,d\log(z_1 - z_2)$.
\item\label{it:drinfeld-face} The Drinfeld $r$-matrix
 \textup{(}for $\cA = \widehat{\fg}_k$\textup{)}:
 $r(z) = \Omega/((k+h^\vee)z)$.
\item The Sklyanin bracket:
 $\{L_1(u), L_2(v)\}
 = [r(u - v), L_1(u) \otimes L_2(v)]$.
\item The Gaudin Hamiltonians:
 $H_i
 = \sum_{j \ne i} r(z_i - z_j)$
 are mutually commuting.
\end{enumerate}
Faces~\textup{(F1)--(F4)} hold for every modular Koszul
algebra. Faces~\textup{(F5)--(F7)} specialize to affine
Kac--Moody; the general case is obtained by replacing
$\Omega/((k{+}h^\vee)z)$ with $r_\cA(z)$.
\end{theorem}

The seven faces are not seven theorems but seven readings of
one equation. The collision residue
$r(z) = \pi_2(\Theta_\cA)$ is the arity-$2$ projection of
the universal MC element; the seven faces are the seven
classical frameworks in which this projection has been
studied. Drinfeld's $r$-matrix~\cite{Drinfeld85}, the KZ
connection~\cite{KZ84}, the Sklyanin
bracket~\cite{STS83}, and the Gaudin model~\cite{FFR94}
all extract the same binary invariant.

The relationship to the modular characteristic is immediate:
\begin{equation}\label{eq:kappa-from-r}
\kappa(\cA)
\;=\;
\mathrm{av}\bigl(r(z)\bigr)\,.
\end{equation}
The averaging map $\mathrm{av}$ is the
$\Sigma_2$-coinvariant projection. For Kac--Moody,
$\mathrm{av}(k\,\Omega/z) = \dim(\fg)(k + h^\vee)/(2h^\vee)$.
For Virasoro, $\mathrm{av}((c/2)/z^3 + 2T/z) = c/2$.

% ====================================================================
\subsection{DS reduction as a functor on collision residues}
\label{ssec:ds-reduction}
% ====================================================================

Drinfeld--Sokolov reduction converts affine Kac--Moody
algebras into $\mathcal{W}$-algebras:
$\widehat{\fg}_k \xmapsto{\mathrm{DS}}
\mathcal{W}_N^k(\fg)$. At the level of the collision
residue, this passage raises the pole order:

\begin{center}
\renewcommand{\arraystretch}{1.2}
\begin{tabular}{lccl}
\toprule
& \textbf{OPE max.\ pole} & \textbf{Collision max.\ pole}
 & \textbf{Class} \\
\midrule
$\widehat{\mathfrak{sl}}_N$ & $2$ & $1$ & $\mathbf{L}$ \\
$\mathcal{W}_N$ & $2N$ & $2N{-}1$ & $\mathbf{M}$ \\
\bottomrule
\end{tabular}
\end{center}

\noindent
The Sugawara construction generates a Virasoro field
$T(z)$ from the affine currents $J^a(z)$. The OPE
$T(z)\,T(w)$ has a quartic pole; DS reduction extends
this mechanism to all generators, producing a
$\mathcal{W}$-algebra whose maximal OPE pole grows with
$N$. The shadow depth jumps from finite (class~$\mathbf{L}$)
to infinite (class~$\mathbf{M}$): DS reduction is the functor
that transports gauge-type algebras into gravitational-type
algebras.

The intertwining theorem controls the Koszul-dual behavior:
the bar-cobar adjunction commutes with principal DS
reduction. For non-principal reductions (hook-type in type
$A$, proved; arbitrary nilpotent, programme), the
intertwining holds on the hook-type corridor and is
conjectural in full generality.


% ====================================================================
\section{The Swiss-cheese realization}
\label{sec:swiss-cheese}
% ====================================================================

The bar complex of a chiral algebra on a curve $X$ has two
structures: a differential (from OPE residues on Fulton--MacPherson
compactifications of~$X$) and a coproduct (from interval-cutting
in a topological direction). The bar complex $B(\cA)$ is an $E_1$
chiral coassociative coalgebra over $(\mathrm{ChirAss})^!$: the
differential encodes holomorphic factorization on~$\mathbb{C}$, the
deconcatenation coproduct encodes topological factorization
on~$\mathbb{R}$. The two-coloured Swiss-cheese operad
$\mathrm{SC}^{\mathrm{ch,top}}$ emerges on the chiral derived center
pair $(C^\bullet_{\mathrm{ch}}(\cA,\cA),\, \cA)$, not on $B(\cA)$
itself.

% ====================================================================
\subsection{The two colours}
\label{ssec:two-colours}
% ====================================================================

The bar complex lives on $\FM_k(\Bbbk) \times \Conf_k(\mathbb{R})$,
the product of holomorphic and topological configuration spaces.
This product is the operadic fingerprint of the Swiss-cheese
operad:

\begin{center}
\renewcommand{\arraystretch}{1.3}
\begin{tabular}{lll}
\toprule
& \textbf{Closed colour} & \textbf{Open colour} \\
\midrule
\textbf{Space} & $\FM_k(\Bbbk)$ & $\Conf_k(\mathbb{R})$ \\
\textbf{Structure} & Bar differential $d_{\barB}$ &
 Deconcatenation $\Delta$ \\
\textbf{Physics} & Holomorphic factorization & Topological
 factorization \\
\textbf{Operadic type} & $E_\infty$ & $E_1$ \\
\textbf{Coalgebra} & $\Sym^c(s^{-1}\bar\cA)$ &
 $T^c(s^{-1}\bar\cA)$ \\
\textbf{Coproduct} & Coshuffle ($2^n$ terms) &
 Deconcatenation ($n+1$ terms) \\
\bottomrule
\end{tabular}
\end{center}

\noindent
The closed colour is the holomorphic factorization of
Section~\ref{sec:bar}: the bar differential extracts OPE
residues, produces $d^2 = 0$ at genus~$0$, and acquires
curvature $\kappa(\cA) \cdot \omega_g$ at higher genus. The
open colour is the topological factorization: the
deconcatenation coproduct splits an ordered sequence at every
consecutive position, producing the cofree tensor coalgebra.

The directionality of the Swiss-cheese operad is strict:
\emph{no open inputs produce closed outputs}. Bulk operators
restrict to boundary operators, but not conversely. The open
colour is the $E_1$-direction; the passage to
$\Sigma_n$-coinvariants (the closed/symmetric bar) is a lossy
projection. The five main theorems A--H are the invariants
that survive this projection.

% ====================================================================
\subsection{Three bar complexes}
\label{ssec:three-bars}
% ====================================================================

Three bar-type objects coexist, corresponding to three
cooperadic structures on the same underlying graded vector
space:

\begin{enumerate}[(a)]
\item $\barB^{\mathrm{ord}}(\cA) = T^c(s^{-1}\bar\cA)$:
 the \emph{ordered bar}. Deconcatenation coproduct ($n + 1$
 terms), coassociative but not cocommutative. Carries the
 $R$-matrix, the KZ associator, and the full Yangian
 deformation. This is the $E_1$-coalgebra.

\item $\barB^{\Sigma}(\cA) = \Sym^c(s^{-1}\bar\cA)$:
 the \emph{symmetric bar}. Coshuffle coproduct ($2^n$
 terms), cocommutative and coassociative. This is the
 factorization coalgebra on $\Ran(X)$: the $E_\infty$-coalgebra
 of Beilinson--Drinfeld.

\item $B^{\mathrm{FG}}(\cA)
 = \mathrm{Lie}^c(s^{-1}\bar\cA)$: the \emph{Francis--Gaitsgory
 bar}. CoLie cobracket, capturing only the zeroth-pole OPE
 data ($a_{(0)}b$, the chiral Lie bracket). A
 quasi-isomorphism $B^{\mathrm{FG}} \hookrightarrow
 \barB^\Sigma$ holds in characteristic zero, but
 $B^{\mathrm{FG}}$ is invisible to the curvature $\kappa(\cA)$:
 the modular characteristic comes from the \emph{first} OPE
 pole, which the Francis--Gaitsgory bar does not see.
\end{enumerate}

The natural descent is
$\barB^{\mathrm{ord}} \xrightarrow{\mathrm{av}}
\barB^{\Sigma} \hookleftarrow B^{\mathrm{FG}}$:
the ordered bar maps to the symmetric bar by
$\Sigma_n$-coinvariance (a lossy projection), and the
Francis--Gaitsgory bar embeds into the symmetric bar
(a quasi-isomorphism in characteristic zero, losing all
higher-pole data). The ordered bar is the natural
primitive; the other two are its images.

% ====================================================================
\subsection{PVA descent and the recognition theorem}
\label{ssec:pva}
% ====================================================================

The cohomology of a Swiss-cheese algebra carries a Poisson
vertex algebra structure: the commutative product descends
from the regular OPE and the Poisson bracket from the singular
OPE. This is the \emph{PVA descent}: the classical shadow of
the full quantum Swiss-cheese structure.

The six PVA axioms (Leibniz, skew-symmetry, Jacobi, sesqui\-linearity,
$\partial$-bracket, and Wick) are proved geometrically from
configuration-space geometry: the exchange cylinder on
$\FM_2(\Bbbk)$ for skew-symmetry, the three-face Stokes
theorem on $\FM_3(\Bbbk)$ for Jacobi, and boundary
factorization on higher $\FM_k(\Bbbk)$ for the remaining
axioms.

The recognition theorem establishes the converse:

\begin{theorem}[Recognition]
\label{thm:recognition}
A factorization algebra on\/ $\Ran(X)$ equipped with Weiss
cosheaf descent data and compatible with the factorization
isomorphisms is a $C_*(\mathcal{W}(\mathrm{SC}^{\mathrm{ch,top}}))$-algebra.
\end{theorem}

The theorem converts the abstract factorization structure
into a concrete operadic algebra over the Boardman--Vogt
resolution of the Swiss-cheese operad. It is proved by Weiss
cosheaf descent: the factorization pattern on $\Ran(X)$
determines the operadic composition uniquely.

The homotopy-Koszulity of $\mathrm{SC}^{\mathrm{ch,top}}$
is proved via Kontsevich formality and transfer from the
classical Swiss-cheese operad (Livernet). This removes all
previously conditional hypotheses: bar-cobar Quillen
equivalence, filtered Koszul duality, and the identification
of line operators with dual modules all become unconditional.


% ====================================================================
\section{Three-dimensional HT QFT and gauge-gravity}
\label{sec:ht-qft}
% ====================================================================

A four-dimensional holomorphic-topological field theory on
$\Bbbk_z \times \mathbb{R}_t \times \mathbb{R}_{>0}$
restricts to a chiral algebra on the holomorphic boundary.
The bar complex of that boundary algebra classifies
twisting morphisms (couplings to the Koszul dual); the
chiral derived centre computes the universal bulk; bar-cobar
inversion recovers the boundary algebra itself. The modular
convolution algebra organizes the full boundary-to-bulk
correspondence at all genera.

% ====================================================================
\subsection{The holographic datum}
\label{ssec:holographic-datum}
% ====================================================================

The complete data of a holographic system is a sextuple:

\begin{definition}[Holographic modular Koszul datum]
\label{def:holographic}
The \emph{holographic modular Koszul datum} of a chiral
algebra~$\cA$ is the sextuple
\begin{equation}\label{eq:holographic-datum}
\mathcal{H}(\cA)
\;=\;
\bigl(\cA,\;
\cA^!,\;
\mathcal{C},\;
r(z),\;
\Theta_\cA,\;
\nabla^{\mathrm{hol}}\bigr)\,,
\end{equation}
where $\cA^! = (H^*(B(\cA)))^\vee$ is the Koszul dual,
$\mathcal{C} \simeq \cA^!\text{-}\mathsf{mod}$ is the
line-operator category, $r(z)$ is the collision residue,
$\Theta_\cA$ is the universal MC element, and
$\nabla^{\mathrm{hol}}_{g,n} = d -
\mathrm{Sh}_{g,n}(\Theta_\cA)$ is the modular shadow
connection. Every component is a projection
of~$\Theta_\cA$.
\end{definition}

The datum $\mathcal{H}(\cA)$ is algorithmically
reconstructed from the boundary OPE of~$\cA$:
$\cA \xmapsto{B} \Theta_\cA \xmapsto{\Res^{\mathrm{coll}}}
r(z) \xmapsto{\mathrm{RTT}} \mathcal{C}$. Each step uses
only the output of the preceding step; the entire holographic
system is determined by the chiral algebra on the boundary.

% ====================================================================
\subsection{The 4d--3d--2d resolution}
\label{ssec:4d-3d-2d}
% ====================================================================

The three volumes of the programme correspond to three levels
of a dimensional hierarchy:

\begin{center}
\renewcommand{\arraystretch}{1.3}
\begin{tabular}{cllll}
\toprule
\textbf{$d$} & \textbf{Level} & \textbf{Operadic type}
 & \textbf{Key datum} & \textbf{Volume} \\
\midrule
$4$ & CY source & $E_2$
 & CY category $\mathcal{C}$ & III \\
$3$ & Swiss-cheese & $E_1$
 & $R$-matrix $r(z)$ & II \\
$2$ & Modular shadow & $E_\infty$
 & $\kappa(\cA)$ & I \\
\bottomrule
\end{tabular}
\end{center}

\noindent
At each level, information is lost. From 4d to 3d, one
topological direction is forgotten: $E_2 \to E_1$ (the
braiding becomes ordering). From 3d to 2d,
$\Sigma_n$-coinvariance projects the $R$-matrix to a scalar:
$r(z) \mapsto \kappa$. The five main theorems are the
invariants that survive both projections.

% ====================================================================
\subsection{Gauge versus gravity}
\label{ssec:gauge-vs-gravity}
% ====================================================================

The shadow depth classification
(Definition~\ref{def:shadow-depth}) separates the standard
landscape into two physical regimes:

\begin{enumerate}[(a)]
\item \emph{Gauge theories} (classes~$\mathbf{G}$,
 $\mathbf{L}$): finite shadow depth, collision residue with
 at most a simple pole.
 The bar complex is effectively quadratic (class~$\mathbf{G}$)
 or cubic (class~$\mathbf{L}$).
 The representation theory is governed by quantum groups with
 rational $R$-matrices.

\item \emph{Gravitational theories} (classes~$\mathbf{C}$,
 $\mathbf{M}$): the shadow tower extends to arity~$4$
 (class~$\mathbf{C}$) or is infinite (class~$\mathbf{M}$).
 The collision residue has higher-order poles. The bar complex
 is genuinely non-formal: the full $A_\infty$-tower is
 irreducible. The Koszul duality
 $\Vir_c^! = \Vir_{26-c}$ (or
 $\mathcal{W}_N^! = \mathcal{W}_N^{k'}$ at the dual level)
 exchanges matter and ghost sectors.
\end{enumerate}

DS reduction is the functor connecting these regimes: it
transports affine Kac--Moody (class~$\mathbf{L}$, gauge) to
$\mathcal{W}$-algebras (class~$\mathbf{M}$, gravity) by
raising the pole order through the Sugawara construction.
The physical content: the Sugawara stress tensor converts
gauge currents into a metric, and the bar complex detects
this conversion as the appearance of higher collision poles.

% ====================================================================
\subsection{The critical string at $c = 26$}
\label{ssec:critical-string}
% ====================================================================

The Virasoro Koszul duality $\Vir_c^! = \Vir_{26 - c}$
identifies three distinguished points:

\begin{center}
\renewcommand{\arraystretch}{1.3}
\begin{tabular}{clll}
\toprule
$\boldsymbol{c}$ & \textbf{Property}
 & $\boldsymbol{\kappa}$ & \textbf{Clifford} \\
\midrule
$0$ & Gauge point: $\kappa = 0$, flat bar
 & $0$ & degenerate \\
$13$ & Self-dual point: $\Vir_{13}^! = \Vir_{13}$
 & $13/2$ & non-degenerate \\
$26$ & Critical string: $\kappa^! = 0$
 & $13$ & exterior algebra \\
\bottomrule
\end{tabular}
\end{center}

\noindent
At $c = 0$: $\kappa(\Vir_0) = 0$, the bar complex is flat at
all genera, and the genus tower is trivial. At $c = 13$:
$\kappa(\Vir_{13}) = 13/2 = \kappa(\Vir_{13}^!)$, the
complementarity pairing is symmetric, and the entire shadow
tower is self-dual. At $c = 26$: $\kappa(\Vir_0^!) = 0$, the
Koszul dual is flat; the Clifford algebra of the
complementarity pairing degenerates to an exterior algebra,
and the matter-ghost cancellation condition
$\kappa_{\mathrm{eff}}
= \kappa(\mathrm{matter}) + \kappa(\mathrm{ghost}) = 0$ of the
bosonic string is satisfied. The $26$-dimensional critical
string is the point where the Koszul dual's bar complex
becomes flat: the gravitational content migrates from curvature
to bar cohomology.


% ====================================================================
\section{The Calabi--Yau bridge}
\label{sec:cy-bridge}
% ====================================================================

Volumes~I and~II take a chiral algebra as given. The source
of chiral algebras is geometry: a Calabi--Yau category
$\mathcal{C}$ of dimension~$d$ produces a chiral algebra
$A_\mathcal{C}$ on a curve~$X$, with the CY trace realized
as the modular characteristic. This is the content of
Volume~III.

% ====================================================================
\subsection{The functor $\Phi$}
\label{ssec:functor-phi}
% ====================================================================

The construction proceeds in three steps. Given a CY$_d$
category $\mathcal{C}$:

\begin{enumerate}[(1)]
\item The cyclic $A_\infty$-structure on $\mathcal{C}$
 determines a \emph{Lie conformal algebra} $L_\mathcal{C}$:
 the cyclic bar differential gives the $\lambda$-bracket,
 the cyclic trace gives the invariant bilinear form, and
 the higher $A_\infty$-operations give higher Lie conformal
 products.

\item The \emph{factorization envelope}
 $U_X^{\mathrm{fact}}(L_\mathcal{C})$ produces a
 factorization algebra on $\Ran(X)$, hence a chiral algebra
 $A_\mathcal{C}$ on~$X$. The PBW theorem for factorization
 envelopes guarantees the expected filtration.

\item The $\mathbb{S}^d$-framing on Hochschild homology
 determines the operadic enhancement. For $d = 2$: the
 $\mathbb{S}^2$-action gives an $E_2$-structure with braided
 monoidal representation category. For $d = 3$: holomorphic
 Chern--Simons theory on $\mathcal{C}$ breaks $E_2$ to $E_1$,
 and the braided structure is recovered via the Drinfeld centre
 of the $E_1$-monoidal category.
\end{enumerate}

\begin{theorem}[CY-A$_2$]
\label{thm:cy-a2}
The functor $\Phi$ is constructed for $d = 2$: all three steps
are proved. The modular characteristic of $A_\mathcal{C}$
equals the Euler characteristic:
$\kappa(A_\mathcal{C}) = \chi^{\mathrm{CY}}(\mathcal{C})$.
\end{theorem}

For $d = 3$, steps~(1) and~(2) are unconditional; step~(3) is
a programme conditional on the chain-level
$\mathbb{S}^3$-framing. The passage from $d = 2$ to $d = 3$
replaces a geometric braiding (from
$\pi_1(\Conf_2(\mathbb{R}^2)) = \mathbb{Z}$) with an
algebraic one (from the Drinfeld centre of the
$E_1$-monoidal representation category): this is the CY
analogue of the Tannakian reconstruction of a quantum group
from its module category.

% ====================================================================
\subsection{The $E_1$ stabilization}
\label{ssec:e1-stabilization}
% ====================================================================

For $d \ge 3$, the CY chiral algebra is $E_1$, not $E_d$.
The reason is structural: the Schouten--Nijenhuis bracket on
polyvector fields of $\Bbbk^d$ vanishes identically for
$d \ge 3$, so the classical Lie conformal algebra is abelian
and all noncommutative structure arises from quantization.
The CoHA construction produces an ordered (associative) product
from the extension correspondence
$0 \to V' \to V \to V'' \to 0$, which has a preferred
direction: sub before quotient. This is $E_1$, not $E_2$.

The framing obstruction lies in $\pi_d(BU)$, which is
$\mathbb{Z}$ for even~$d$ and $0$ for $d \equiv 1, 3, 7
\pmod{8}$. The pattern repeats with period~$8$ by Bott
periodicity. For $d = 2$: the obstruction is
$c_1 \in \pi_2(BU) = \mathbb{Z}$, the braiding parameter.
For $d = 3$: the obstruction vanishes
($\pi_3(BU) = 0$), so the $\mathbb{S}^3$-framing trivializes
and the $E_1$-structure receives no additional enhancement.
For $d = 4$: $\pi_4(BU) = \mathbb{Z}$ produces a shifted
symplectic level ($p_1$), the CY$_4$ analogue of the
Kac--Moody level.

% ====================================================================
\subsection{The $K3 \times E$ prototype}
\label{ssec:k3-times-e}
% ====================================================================

The product of a K3 surface~$S$ and an elliptic curve~$E$
is the canonical CY$_3$ testing ground: simple enough to
compute, rich enough to obstruct.

The lattice $\Lambda^{3,2}$ has signature $(3,2)$; the
Jacobi form $\phi_{0,1}(\tau, z)$ is the K3 elliptic genus;
the Igusa cusp form $\Delta_5$ is the Borcherds denominator
identity. The modular characteristic of the chiral de Rham
complex is $\kappa_{\mathrm{ch}}(K3 \times E) = 3
= \dim_\Bbbk$, verified by six independent paths. The BKM
automorphic weight is a distinct quantity:
$\kappa_{\mathrm{BKM}} = 5$ (the weight of $\Delta_5$).

These are different numbers measuring different things.
A single CY$_3$ admits multiple chiral algebraizations,
each producing a different modular characteristic. The
$\kappa$-spectrum
\begin{equation}\label{eq:kappa-spectrum}
\operatorname{Spec}_\kappa(K3 \times E)
\;=\; \{2, 3, 5, 24\}
\end{equation}
records four values: $\kappa_{\mathrm{cat}} = 2
= \chi(\mathcal{O}_{K3})$ (the holomorphic Euler
characteristic), $\kappa_{\mathrm{ch}} = 3$ (the complex
dimension), $\kappa_{\mathrm{BKM}} = 5$ (the Borcherds
weight), and $\kappa_{\mathrm{fiber}} = 24$ (the lattice
rank). Four algebras see four aspects of the same manifold.

Four structural obstructions separate the shadow obstruction
tower from the full automorphic form:
\begin{enumerate}[label=\textup{(O\arabic*)}]
\item \emph{Categorical}: the shadow tower produces Taylor
 coefficients; the denominator identity is an automorphic form
 on $\mathcal{H}_2$.
\item \emph{$\kappa$-mismatch}:
 $\kappa_{\mathrm{ch}} = 3 \ne 5 = \kappa_{\mathrm{BKM}}$.
\item \emph{Second quantization}: the physical DT partition
 function is the DMVV symmetric product of single-copy data;
 the bar complex is the single copy.
\item \emph{Schottky}: at genus $g \ge 4$, the Torelli map
 $\cM_g \to \mathcal{A}_g$ is no longer dominant, and the
 shadow tower is imprisoned in tautological classes.
\end{enumerate}
Each obstruction points to a research programme. The
$\kappa$-spectrum is the first invariant that separates CY
chiral algebras with the same underlying manifold.


% ====================================================================
\subsection{$E_1$ at $d = 3$ via holomorphic Chern--Simons}
\label{ssec:e1-holcs}
% ====================================================================

For a CY threefold $X$, holomorphic Chern--Simons theory
on $X$ with gauge group~$G$ produces a chiral algebra
$A_X$ of $E_1$-type. The spectral parameter of the
$R$-matrix is identified with the holomorphic coordinate
of the 4d theory restricted to a 2d boundary. The
braided monoidal structure on the representation category is
recovered through the Drinfeld centre:
\begin{equation}\label{eq:drinfeld-centre}
\mathcal{Z}\bigl(\mathrm{Rep}^{E_1}(A_X)\bigr)
\;\simeq\;
\mathrm{Rep}^{E_2}(A_X)\,,
\end{equation}
where the right-hand side is the braided monoidal category of
$E_2$-modules. For toric CY$_3$, the $E_1$-algebra is the
critical CoHA $\mathrm{CoHA}(Q_X)$: the positive half of an
affine super Yangian $Y(\widehat{\fg}_{Q_X})$. For
$\Bbbk^3$: the Jordan quiver gives
$Y(\widehat{\mathfrak{gl}}_1) \simeq \mathcal{W}_{1+\infty}$.

The dimensional reduction principle underlies the full trilogy:
$4d \to 3d$ forgets the braiding ($E_2 \to E_1$); $3d \to 2d$
forgets the ordering ($E_1 \to E_\infty$ by
$\Sigma_n$-coinvariance). The five main theorems live at the
$2d$ level; the seven faces of the collision residue live at
the $3d$ level; the CY source lives at the $4d$ level. The
modular Koszul engine converts enumerative geometry
(Donaldson--Thomas invariants, root multiplicities, denominator
identities) into homotopy algebra (bar complexes, MC elements,
genus expansions). One Maurer--Cartan element
$\Theta_{A_\mathcal{C}}$ absorbs the entire automorphic datum.


% ====================================================================
% Additional references for Sections 5--8
% ====================================================================

% NOTE: When merging into the main document, add these entries
% to the thebibliography environment of Section 1.
% They are listed here as comments for reference; the active
% bibliography is in programme_summary_sections9_14.tex.

% \bibitem{Drinfeld85}
% V.\,G.~Drinfeld,
% Hopf algebras and the quantum Yang--Baxter equation,
% \textit{Dokl.\ Akad.\ Nauk SSSR}~\textbf{283} (1985), 1060--1064.

% \bibitem{FFR94}
% B.~Feigin, E.~Frenkel, and N.~Reshetikhin,
% Gaudin model, Bethe ansatz and critical level,
% \textit{Comm.\ Math.\ Phys.}~\textbf{166} (1994), 27--62.

% \bibitem{KZ84}
% V.\,G.~Knizhnik and A.\,B.~Zamolodchikov,
% Current algebra and Wess--Zumino model in two dimensions,
% \textit{Nuclear Phys.\ B}~\textbf{247} (1984), 83--103.

% \bibitem{STS83}
% M.\,A.~Semenov-Tian-Shansky,
% What is a classical $r$-matrix?,
% \textit{Funct.\ Anal.\ Appl.}~\textbf{17} (1983), 259--272.

% \bibitem{DNP25}
% T.~Dimofte, N.~Niu, and A.~Py,
% Twisted holography and Koszul duality,
% preprint, 2025.

% \bibitem{CWY18}
% K.~Costello, E.~Witten, and M.~Yamazaki,
% Gauge theory and integrability, I--II,
% \textit{ICCM Not.}~\textbf{6} (2018), 46--119 and 120--149.
